%A brief introduction to the project problem you are addressing, contextualising the problem, motivating it, and explaining its importance.
The Traveling Salesman Problem (TSP) is one of the most extensively studied problems in combinatorial optimization. Given a set of cities and pairwise travel costs, the objective is to determine a minimum-cost Hamiltonian cycle visiting each city exactly once before returning to the starting point. The TSP is known to be NP-hard, as it can be related through polynomial reductions to the Hamiltonian Cycle problem.

Beyond its theoretical importance, the TSP constitutes the foundation of numerous routing and scheduling problems in operations research. Consequently, many variants have emerged in the literature to model additional real-world constraints such as capacities, profits, precedence relations, and time windows.

The Orienteering Problem (OP) is a variant of the TSP in which each vertex is associated with a profit or score. Unlike the classical TSP, the OP seeks to determine a subset of vertices maximizing the collected score while satisfying a constraint on the maximum allowable tour length. The problem was first introduced by Tsiligirides, where it was referred to as the Generalized Traveling Salesman Problem (GTSP).

In this project, the Orienteering Problem with Time Windows (OPTW) is considered, in which a time interval $[O_i, C_i]$ is associated with each vertex. The objective is identical to the OP, with the additional constraint that each vertex must either be visited within its time window or skipped entirely.

This project applies the OPTW to the practical problem of touristic landmark visit planning in the city of Algiers. Given the density and diversity of Algiers' cultural and historical sites --- many of which have restricted visiting hours --- the OPTW provides a natural and expressive model for planning a day's itinerary that maximises the interest of visited landmarks while respecting opening times and the tourist's available time budget. The system takes as input a set of landmarks with associated interest scores, opening hours, and visit durations, and produces an ordered daily route starting and ending at the tourist's hotel. A demonstration interface was developed to allow users to specify their starting time, available duration, and travel day, and to visualise the resulting optimised route on a map of Algiers. The practical benefit to the end user is a ready-to-follow itinerary that avoids wasted travel and closed sites, making the most of a limited visit to the city.